\documentclass[10pt,a4paper]{amsart}
\usepackage[margin=19mm]{geometry}
\usepackage{amsmath,amssymb,mathtools}
\usepackage[scr=rsfs,cal=euler]{mathalfa}
\usepackage{xfrac}
\usepackage[unicode,hidelinks,bookmarks=false]{hyperref}
\newcommand{\ie}{i.e.\ }
\newcommand{\cf}{cf.\ }
\newcommand{\resp}{respectively\ }
\newcommand{\Subsection}{\subsection}
\providecommand{\nobreakdash}{\nobreak\hbox{-}}
\setlength{\emergencystretch}{2em}
\multlinegap0pt
\usepackage{bm}
\usepackage{bbm}
\usepackage{dsfont}
\usepackage{xspace}
\usepackage{cancel}
\usepackage{mathtools}
\usepackage[shortlabels]{enumitem}
\usepackage{amsthm}
\makeatletter
\@ifundefined{theorem}{\newtheorem{theorem}{Theorem}}{}
\makeatother
\DeclareMathOperator{\od}{ObsDiam}
\DeclareMathOperator{\pd}{PartDiam}
\newcommand{\R}{\mathbb{R}}
\title[Extensions of One-Sided Box Geometry and Pyramid Invariants ]{Extensions of One-Sided Box Geometry and Pyramid Invariants to gd-Sets and qm-Spaces}
\author{Shigeaki Yokota}
\date{Source: arXiv:2608.12749v1 (CC BY 4.0)}
\begin{document}
\maketitle

\noindent\emph{Source and licence.} Extensions of One-Sided Box Geometry and Pyramid Invariants to gd-Sets and qm-Spaces. Version arXiv:2608.12749v1 (CC BY 4.0), distributed under the Creative Commons Attribution 4.0 International licence, as recorded in the arXiv licence metadata for this exact version and shown on the abstract page https://arxiv.org/abs/2608.12749v1. Full rights notice: licenses/arxiv-2608.12749-CC-BY-4.0.txt.\par\smallskip

\noindent\textit{Editorial scope.} Counted unit: the Theorem whose source label is \texttt{thm:elementary-feature-concentration-comparison} in the article (source0.tex lines 1512--1543, source label \texttt{thm:elementary-feature-concentration-comparison}), delivered together with the article's own complete original proof of it (source0.tex lines 1545--1627, one \texttt{proof} environment of 83 lines). No mathematics is written, repaired or re-derived: statement and proof are the article's own text line for line. Required context retained uncounted: none. Every definition, display and label of those lines is retained unchanged. Omitted from this package and not redistributed: 1--1511, 1544--1544, 1628--2130. Nothing inside a retained excerpt is omitted or altered: every retained body is delivered exactly as the article's lines are written, so no omission receipt of any kind accompanies this package. Citations: the retained excerpts carry no citation command at all, so no bibliography entry of the article is transcribed and no reference is redistributed. Reference adaptation: none was needed and none was made; every reference inside the delivered unit resolves inside it, so no unresolved reference or citation is produced. Printed slips kept literal: no printed word, symbol, hypothesis or number of the counted statement or of its proof was altered. Layout and class: the article's own class is \texttt{amsart} with packages including amsmath, amssymb, amsthm, mathtools, enumitem, iftex, hyperref, cleveref; the wrapper is the lane's pinned \texttt{amsart} editorial wrapper at 10pt on A4, which restates the theorem-like environments the retained text needs and adds the article's own macro definitions that the retained text needs (\texttt{\textbackslash R}, \texttt{\textbackslash od}, \texttt{\textbackslash pd}), with extra wrapper packages (bm, bbm, dsfont, xspace, cancel, mathtools, enumitem). The delivered numbering is the unit's own. Assets: no retained span contains a figure environment, a graphics-inclusion command, a PDF-page inclusion, a table or a TikZ picture; no image file is copied into this package, the package declares no asset dependency and no image file is redistributed with it. Licence: version arXiv:2608.12749v1 of this article is distributed under the Creative Commons Attribution 4.0 International licence, as recorded in the arXiv licence metadata for this exact version and shown on the abstract page https://arxiv.org/abs/2608.12749v1. A full rights notice is delivered at licenses/arxiv-2608.12749-CC-BY-4.0.txt. Characters: every retained excerpt is pure ASCII, so the delivered engine, XeLaTeX with its default Latin Modern text font, reports no missing character.
\begin{theorem}[Elementary comparisons for function-family concentration]
\label{thm:elementary-feature-concentration-comparison}
Let $X$ be a gd-set.  Then the following hold.
\begin{enumerate}[label=\textup{(\roman*)}]
\item If $0<\kappa<1/2$, then
\[
\od(X;-2\kappa)\leq\operatorname{Sep}^+(X;\kappa,\kappa).
\]
\item If $0<\kappa'<\kappa\leq1/2$, then
\[
\operatorname{Sep}^+(X;\kappa,\kappa)\leq\od(X;-\kappa').
\]
\item If $0<\kappa<1$ and $r_+,r_->0$ satisfy
\[
\alpha_X^+(r_+)+\alpha_X^-(r_-)\leq\kappa,
\]
then
\[
\od(X;-\kappa)\leq r_++r_-.
\]
\item For $r>0$,
\[
\max\{\alpha_X^+(r),\alpha_X^-(r)\}
\leq\sup\{\kappa\in(0,1)\mid\od(X;-\kappa)\geq r\}.
\]
If the set on the right is empty, its supremum is understood to be $0$.
\end{enumerate}
These comparisons require no closure of the function family under maximum or
minimum, truncation, postcomposition by $1$-Lipschitz maps, or
inf-convolution.  They also require neither the inclusion of constant
functions nor $0\in\overline{F_X}$.
\end{theorem}

\begin{proof}
For \textup{(i)}, use the lower and upper quantile sets.  Take
$f\in\overline{F_X}$ and put $\nu\coloneqq f_*\mu_X$.  Define
\[
\rho_-\coloneqq\sup\{t\in\R\mid\nu((-\infty,t))\leq\kappa\},\qquad
\rho_+\coloneqq\inf\{t\in\R\mid\nu((t,+\infty))\leq\kappa\}.
\]
The elementary properties of quantiles give
\[
\nu((-\infty,\rho_-])\geq\kappa,\quad
\nu([\rho_+,+\infty))\geq\kappa,\quad
\nu([\rho_-,\rho_+])\geq1-2\kappa,\quad
\rho_-\leq\rho_+.
\]
Therefore,
$\pd(\nu;1-2\kappa)\leq\rho_+-\rho_-$.  If
$\rho_-<\rho_+$, then
\[
A_0\coloneqq f^{-1}((-\infty,\rho_-]),\qquad
A_1\coloneqq f^{-1}([\rho_+,+\infty))
\]
are disjoint, each has measure at least $\kappa$, and
\[
\sigma^\to(X;A_0,A_1)
\geq\inf_{a\in A_0,\,b\in A_1}(f(b)-f(a))
\geq\rho_+-\rho_-.
\]
If $\rho_-=\rho_+$, the partial diameter is $0$.  Thus, in either case,
\[
\pd(f_*\mu_X;1-2\kappa)
\leq\operatorname{Sep}^+(X;\kappa,\kappa).
\]
Taking the supremum over $f$ proves \textup{(i)}.

To prove \textup{(ii)}, take disjoint Borel sets $A_0,A_1$ such that
\[
\mu_X(A_i)\geq\kappa\qquad(i=0,1).
\]
If
$\sigma^\to(X;A_0,A_1)\leq0$, the assertion follows from nonnegativity.
Otherwise, take $0<s<\sigma^\to(X;A_0,A_1)$ and choose
$f\in\overline{F_X}$ such that
\[
\inf_{a\in A_0,\,b\in A_1}(f(b)-f(a))>s.
\]
Let $C\subset\R$ be a Borel set such that
$f_*\mu_X(C)\geq1-\kappa'$.  For $i=0,1$,
\[
\mu_X(f^{-1}(C))+\mu_X(A_i)\geq1-\kappa'+\kappa>1,
\]
and hence $f^{-1}(C)\cap A_i\neq\varnothing$.  Thus $C$ meets both
$f(A_0)$ and $f(A_1)$.  Since $\operatorname{diam}C>s$, we have
\[
\pd(f_*\mu_X;1-\kappa')\geq s.
\]
Let $s\uparrow\sigma^\to(X;A_0,A_1)$ and take the supremum over
$A_0,A_1$ to obtain \textup{(ii)}.

For \textup{(iii)}, take $f\in\overline{F_X}$ and
$m\in\operatorname{Med}(f)$.  The assumption gives
\[
\mu_X(m-r_-<f<m+r_+)
\geq1-\alpha_X^-(r_-)-\alpha_X^+(r_+)\geq1-\kappa.
\]
The interval has diameter $r_++r_-$, so
$\pd(f_*\mu_X;1-\kappa)\leq r_++r_-$.  Taking the supremum over $f$ proves
the assertion.

Consider the upper tail in \textup{(iv)}.  Take
$f\in\overline{F_X}$, $m\in\operatorname{Med}(f)$, and
$0<\kappa<\mu_X(f\geq m+r)$.  The set $\{f\leq m\}$ has measure at least
$1/2$ and is disjoint from the upper tail, so $\kappa<1/2$.  Let
$C\subset\R$ be any Borel set such that
$f_*\mu_X(C)\geq1-\kappa$.  Since
$\mu_X(f^{-1}(C))\geq1-\kappa$ and
$\mu_X(f\leq m)\geq1/2>\kappa$, the set $C$ meets
$(-\infty,m]$.  It also meets $[m+r,+\infty)$ because
$\mu_X(f\geq m+r)>\kappa$.  Therefore,
$\operatorname{diam}C\geq r$, and hence $\od(X;-\kappa)\geq r$.  Letting
$\kappa\uparrow\mu_X(f\geq m+r)$ and taking the supremum over $f$ and $m$
give the estimate for $\alpha_X^+(r)$.  For $\alpha_X^-(r)$, use
$[m,+\infty)$ and $(-\infty,m-r]$.  This completes the proof.
\end{proof}

\end{document}
